% Generated draft: requires mathematical and source review.
\documentclass[11pt]{article}
\usepackage[margin=25mm]{geometry}
\usepackage{fontspec}
\setmainfont{Latin Modern Roman}
\usepackage{amsmath,amssymb,mathtools,unicode-math}
\setmathfont{Latin Modern Math}
\usepackage{enumitem,needspace,xurl,hyperref}
\setlength{\parindent}{0pt}
\setlength{\parskip}{5pt}
\setlength{\emergencystretch}{3em}
\newcommand{\N}{\mathbb{N}}
\newcommand{\Z}{\mathbb{Z}}
\newcommand{\Q}{\mathbb{Q}}
\newcommand{\R}{\mathbb{R}}
\newcommand{\C}{\mathbb{C}}
\newcommand{\F}{\mathbb{F}}
\newcommand{\A}{\mathbb{A}}
\newcommand{\PP}{\mathbb{P}}
\newcommand{\E}{\mathbb{E}}
\newcommand{\eps}{\varepsilon}
\newcommand{\dd}{\,\mathrm{d}}
\begin{document}
\section*{UnsolvedMath ID: 30006256}
\textbf{Source problem code:} catalogue entry
\subsection{Definitions and mathematical statement}
Let \(X=G/P\) denote a flag variety in the class intended by the supplied record, and let \(T\) denote its equivariance torus. The record does not specify the ground field, the admissible groups \(G\), the parabolic subgroups \(P\), or the relationship between \(T\) and \(G\). Write \(QK_T(X)\) for the \(T\)-equivariant quantum \(K\)-theory ring named in the record, and write \(\star\) for its quantum product. For an algebraic torus \(T\), its representation ring is \(R(T)=\Z[X^*(T)]\), where \(X^*(T)\) is its character lattice.

To express the requested property schematically, let \(\{b_u\}_{u\in I}\) be the intended Schubert basis, let \(D\) be the intended monoid of effective curve degrees, and let \(q^d\) be the formal Novikov monomial corresponding to \(d\in D\). The product coefficients are denoted by
\[
 b_u\star b_v=\sum_{w\in I}\sum_{d\in D}c_{u,v}^{w,d}\,q^d b_w.
\]
This notation is schematic: the basis convention, coefficient extension or localization, and any required completion are not specified by the record.

\textbf{Requested conjecture, with unresolved parameters.} For every flag variety in the intended class, do all coefficients \(c_{u,v}^{w,d}\), for all \(u,v,w\in I\) and \(d\in D\), satisfy the intended alternating-sign rule? A precise assertion of this kind would require a specified parity function \(\varepsilon(u,v,w,d)\in\Z/2\Z\) and a specified positive cone \(C_T\) in the chosen coefficient ring, so that the condition could be written
\[
 (-1)^{\varepsilon(u,v,w,d)}c_{u,v}^{w,d}\in C_T
 \qquad(u,v,w\in I,\ d\in D).
\]
Neither \(\varepsilon\) nor \(C_T\) is supplied. They are placeholders for missing data, not definitions of the intended sign rule. In particular, this display does not specify a complete mathematical conjecture.

The supplied statement also describes an equivalent alternating-sign assertion for a Pontryagin product in localized equivariant \(K\)-theory of an affine Grassmannian. It does not specify the affine Grassmannian, the precise \(K\)-theory variant, the localization, the basis, the sign rule, or the correspondence with \(QK_T(G/P)\). This claimed equivalence is therefore retained as a source qualification and cannot be formulated precisely from the supplied evidence.
\subsection{Short English statement}
The supplied problem asks whether Schubert product coefficients in equivariant quantum \(K\)-theory of flag varieties obey an alternating-sign positivity rule, and asserts a corresponding Pontryagin-product formulation. The actual sign formula, positivity convention, geometric scope, and comparison map are missing, so this draft records a conjecture schema requiring mathematical review.
\subsection{Sources}
\begin{itemize}
\item Supplied local UnsolvedMath registry record: permanent ID \(30006256\), title \emph{Alternating Signs in Equivariant Quantum K-Theory}, source problem code \texttt{catalogue entry}; Oberwolfach Reports Open Problems collection, proposed year \(2025\).
\item Supplied report citation: \emph{Algebraic Groups} (2025), Oberwolfach Reports, pp.\ 941--1010, DOI \texttt{10.4171/owr/2025/20}. Supplied source URL: \url{https://doi.org/10.4171/owr/2025/20}. No problem-specific page or conjecture number is supplied.
\item Supplied literature source, labeled ``OWR 2025 report'': \url{https://ems.press/content/serial-article-files/51849}.
\item Supplied literature source, labeled ``Conjectural positivity preprint'': \url{https://borowiecki.dev/pdf/2510.07689}. The record supplies no authors, full title, or conjecture number for this reference.
\item Attribution evidence: the original extracted text places Shrawan Kumar immediately after the phrase describing alternating-sign behavior. It then contains the title \emph{Center of the small quantum group: towards a combinatorial model}, followed by Anna Lachowska and a joint-work credit to Qi You, Nicolas Hemelsoet, and Oskar Kivinen. The extraction does not establish the boundary between these contributions; its names cannot all be assigned to the present conjecture. The structured proposer field is empty.
\item Literature assessment dated exactly \texttt{2026-08-22}: the supplied curation labels the problem open and states that the alternating-sign conjecture for the equivariant quantum \(K\)-theory Pontryagin product is known only in special cases, with the general conjecture remaining open. This is supplied dated curation, not an independently verified literature assessment. The background expressly requests independent verification of time-sensitive status and cited claims.
\item The registry labels the statement \texttt{corrected\_verified} and reports that the official Oberwolfach report was checked, the question isolated, and notation repaired. These are supplied provenance claims; no source consultation or mathematical verification is asserted by this draft.
\end{itemize}
\end{document}
